\section{The full gamma saddle and static inner overlap}\label{sec:full-saddle}
Keeping the full implicit gamma saddle and its Gaussian determinant
removes the exponential time restriction of the linearized model.
The gain is a joint static estimate, uniform at every real height for
sufficiently large positive time, together with actual spatial overlap
with the inner clock. The field still requires a specified
special-function saddle and the fixed spatial average. We consolidate
Note~30 \cite{n30}; all thresholds below are eventual.

Fix $\ell,h>0$, $Y\ge0$ and the kernel of \eqref{eq:kernel}.
The holomorphic domains include a fixed enlargement of
\[
 |\Re z-x|\le\ell+2,\qquad |\Im z|\le Y+h+2,
\]
large enough to contain the full disks of every root collar meeting
the averaging support. All constants may depend on these fixed data
and the regulator parameter, but are independent of $t,x$ in their
declared domains. Write
\begin{equation}
 X=|x|,\quad \mathcal R=2+t+X,\quad \mathcal L=\log\mathcal R,
 \quad N_*=1+X+t\mathcal L,\qquad e_*=\frac{t^{3/2}}{N_*^2}.
 \label{eq:full-saddle-scales}
\end{equation}
The exact gamma factor $\mathcal M$ and $\alpha=\mathcal M'/\mathcal M$
are those of Section~\ref{sec:gamma-outer}; in particular
$\mathcal M(w)=w(w-1)\pi^{-w/2}\Gamma(w/2)/16$, distinct from
the real source mass $M_0$ and from a Stirling model.

For $w=s=(1+iz)/2$ or $1-s$, define
\begin{align}
 \mathfrak p_t(w)&=w+\frac t2\alpha(\mathfrak p_t(w)),&
 v_*&=\mathfrak p_t(w)-w,&
 B&=1-\frac t2\alpha'(\mathfrak p_t(w)),\notag\\
 \mathcal S_t(w)&=\frac{\mathcal M(\mathfrak p_t(w))e^{-v_*^2/t}}
                         {\sqrt B},&
 P_*(t,z)&=\mathcal S_t(s)+\mathcal S_t(1-s).
 \label{eq:full-saddle-model}
\end{align}
Here $\mathfrak p_t$ is the large branch specified below and $\sqrt B$
is near $1$. This defines a holomorphic field on the required buffered
domains, and on the whole fixed-width horizontal strip at late time;
no entire continuation through all possible saddle branches is asserted.
Use the physical quantities $S_F,w_F,U_F,W_F,\rho_F,R_F$ from
Section~\ref{sec:gamma-outer}, always with the same $q=\epsilon_g^2$
for $F=H,P_*$.

\begin{theorem}[Joint full-saddle static comparison]\label{thm:full-saddle-static}
There are $C,R_*>0$, depending only on the fixed buffers, kernel and
regulator parameter, such that, for $t\ge32$, $\mathcal R\ge R_*$ and
$|y|\le Y$,
\begin{equation}
 |R_H-R_{P_*}|+
 \sum_{|\gamma|\le4}\left(
 |\partial^\gamma(U_H-U_{P_*})|
 +|\partial^\gamma(W_H-W_{P_*})|\right)
 \le C\frac{t^{3/2}}{(1+X+t\log(2+t+X))^2}.
 \label{eq:full-saddle-static}
\end{equation}
The multi-index $\gamma$ is spatial. On the enlarged holomorphic domain,
\begin{equation}
 |H-P_*|\le Ce_*\bigl(|\mathcal S_t(s)|+|\mathcal S_t(1-s)|\bigr).
 \label{eq:full-saddle-field}
\end{equation}
The denominator implicit in this field comparison is the sum of carrier
amplitudes, which stays positive, rather than $|P_*|$ at its zeros.
\end{theorem}

\subsection{The branch and the complete field estimate}
\begin{lemma}[Specified large saddle branch]\label{lem:full-saddle-branch}
Let $w=\sigma+i\tau_0$, with $\sigma$ in the fixed buffered interval.
For $\tau_0\ge0$ and sufficiently large $\mathcal R$, there is exactly
one solution of \eqref{eq:full-saddle-model} in the rectangle
\begin{equation}
 \sigma+t\mathcal L/8\le\Re\mathfrak p\le\sigma+t\mathcal L,
 \qquad \tau_0\le\Im\mathfrak p\le\tau_0+t.
 \label{eq:full-saddle-rectangle}
\end{equation}
Replacing $X$ by $2|\tau_0|$ only changes the scales by fixed constants.
The reflected rectangle specifies negative $\tau_0$. On these branches,
\begin{equation}
 a:=\Re v_*\asymp t\mathcal L,\quad
 |h_*|:=|\Im v_*|\le t,\quad
 |\mathfrak p_t(w)|\asymp N_*,\quad |B-1|\le C/\mathcal L,
 \qquad a\ge t\mathcal L/5
 \label{eq:full-saddle-branch-bounds}
\end{equation}
after increasing the onset. The branch is holomorphic in $w$ on
overlapping buffered domains.
\end{lemma}
\begin{proof}
Sectorial gamma asymptotics and Cauchy differentiation of their
holomorphic remainder on a slightly larger fixed sector give
\begin{equation}
 \alpha(\mathfrak p)=\tfrac12\operatorname{Log}(\mathfrak p/(2\pi))
                         +O(|\mathfrak p|^{-1}),\qquad
 \alpha^{(j)}(\mathfrak p)=O_j(|\mathfrak p|^{-j}),\quad j\ge1.
 \label{eq:full-saddle-sector}
\end{equation}
All saddles stay in the right half-plane, with a fixed sector margin
(for example inside $|\arg\mathfrak p|\le3\pi/4$).
The corresponding bounds on horizontal rays with real part at least
$2$ extend to bounded modulus by increasing constants. These are the
gamma estimates of \cite{DLMFGammaOuter2026}, rather than a theta
large-time theorem. Throughout \eqref{eq:full-saddle-rectangle},
$|\mathfrak p|\asymp N_*$ and
$\log|\mathfrak p|=\mathcal L+O(\log\mathcal L)$.
The fixed-point map has real increment
$t\mathcal L/4+O(t\log\mathcal L)$.

To check its nonnegative imaginary increment, put $A=\Re\mathfrak p$.
For $0\le\tau\le A$ the real part of
$\alpha'(A+i\tau)=1/[2(A+i\tau)]+O(A^{-2})$ is positive.
Vertical integration from the real axis gives $\Im\alpha\ge0$.
For $\tau\ge A$, its leading half-argument is at least $\pi/8$,
and the remainder is smaller; also $\Im\alpha\le1$ eventually.
The map therefore sends the rectangle into itself. Its derivative is
bounded by $Ct/N_*\le C/\mathcal L<1/2$, so contraction on this convex
rectangle proves existence and uniqueness there. It does not prove
uniqueness among every solution of the meromorphic saddle equation.
The more precise identity
$a=(t/4)\log(|\mathfrak p|/(2\pi))+O(t/|\mathfrak p|)$
gives \eqref{eq:full-saddle-branch-bounds}.
Conjugation gives the negative branch, including a real saddle when
$\tau_0=0$. The implicit function theorem and $B\ne0$ give holomorphic
extension across the real axis; uniqueness makes overlapping branches
agree. The logarithm of $\mathcal M$ in $\Re\mathfrak p>1$ is real on
the positive real axis, and $B$ stays in a fixed disk about $1$.
\end{proof}

\begin{lemma}[Both arms, all modes and the middle]\label{lem:full-saddle-field}
The field bound \eqref{eq:full-saddle-field} holds before making any
comparison involving the regulator.
\end{lemma}
\begin{proof}
Start with the exact complete convolution
\begin{equation}
 H(t,z)=\frac1{\sqrt{\pi t}}\int_\R e^{-v^2/t}\frac{\xi(s+v)}8\,dv.
 \label{eq:full-saddle-convolution}
\end{equation}
It is the canonical identity also appearing as equation~(35) of
\cite{polymath2019heat}; its validity here does not use that paper's
$0<t\le1/2$ high-zero theorem. Split at real parts $2$ and $-1$.
Exact reflection $\xi(w)=\xi(1-w)$ and $v\mapsto-v$ turn the two
unbounded arms into
\begin{equation}
 I_w=\frac1{\sqrt{\pi t}}\int_{v_0}^\infty
   e^{-v^2/t}\mathcal M(w+v)\zeta(w+v)\,dv,
 \qquad v_0=2-\Re w,\qquad w=s,1-s.
 \label{eq:full-saddle-arms}
\end{equation}
The remaining middle has bounded length. This pointwise decomposition
is used for estimates; the complete field and the model are holomorphic
independently of the moving endpoints.

Move each arm to $v=r+ih_*$, $r\ge v_0$, and retain the oriented
connector from $v_0$ to $v_0+ih_*$. The original arm is the shifted
horizontal integral plus that connector. First close a finite rectangle
and then send its right endpoint to infinity. Its right face vanishes,
since for fixed parameters the integrand is bounded by
$\exp(-r^2/t+O(r\log(r+2)))$.
Every arm and connector has $\Re(w+v)\ge2$, so no gamma singularity
or zeta pole is crossed and the ordinary full Dirichlet series is valid.
The same oriented construction works for $h_*<0$.

Put $\mathfrak p=\mathfrak p_t(w)$, $m=|\mathfrak p|$ and
$v=v_*+u$, $u\in\R$. For $u\ge-a/2$ define
\begin{align}
 \Gamma_{\mathfrak p}(u)
 &=\log\mathcal M(\mathfrak p+u)-\log\mathcal M(\mathfrak p)
                                       -\alpha(\mathfrak p)u,\notag\\
 E_{\mathfrak p}(u)&=\Gamma_{\mathfrak p}(u)
                              -\tfrac12\alpha'(\mathfrak p)u^2.
 \label{eq:full-saddle-remainders}
\end{align}
The whole intervening segment lies in $\Re>1$ with modulus at least
$cm$; on the positive ray that modulus only increases. Thus Taylor's
integral remainder and \eqref{eq:full-saddle-sector} give on this
\emph{whole} half-line
\[
 |\Gamma_{\mathfrak p}(u)|\le Cu^2/m,
 \qquad |E_{\mathfrak p}(u)|\le C|u|^3/m^2.
\]
The exact first-order cancellation and the retained full quadratic term
are
\begin{equation}
 \frac{e^{-(v_*+u)^2/t}\mathcal M(\mathfrak p+u)}
      {e^{-v_*^2/t}\mathcal M(\mathfrak p)}
       =e^{-Bu^2/t+E_{\mathfrak p}(u)}.
 \label{eq:full-saddle-cancellation}
\end{equation}
Since $t/m\le C/\mathcal L$, both the exact exponential modulus in
\eqref{eq:full-saddle-cancellation} and $|e^{-Bu^2/t}|$ are at most
$e^{-cu^2/t}$, with a fixed $c>0$ (one may take $c=1/2$).
Using $|e^E-1|\le|E|\max(1,e^{\Re E})$, rather than requiring a
small cubic remainder at every $u$, yields the decisive payment
\begin{equation}
 \frac1{\sqrt{\pi t}}\int_{-a/2}^\infty
  |e^{-Bu^2/t}(e^{E_{\mathfrak p}(u)}-1)|\,du
 \le C\frac{t^{3/2}}{m^2}.
 \label{eq:full-saddle-cubic-payment}
\end{equation}
The full complex Gaussian is $B^{-1/2}$ because $\Re B>0$.
Its missing left tail and the true right tail beyond $u=a/2$ have
bound $Ce^{-ca^2/t}$. There is no $e^{t(\Im\alpha)^2/4}$ loss.

For the remaining horizontal segment $v_0\le r\le a/2$, put
$d=a-r$, $A=\Re\mathfrak p$, $\tau=\Im\mathfrak p$.
The inequality $|v+i\tau|/|A+i\tau|\ge v/A$, $2\le v\le A$,
and \eqref{eq:full-saddle-sector} imply
\begin{equation}
 \int_{A-d}^A\Re\alpha(v+i\tau)\,dv
 \ge d\Re\alpha(\mathfrak p)-CA.
 \label{eq:full-saddle-left-integral}
\end{equation}
Indeed the possible logarithmic loss is at most
$-(1/2)\int_0^A\log(v/A)\,dv=A/2$, and the bounded remainders
contribute $O(A)$. Since $\Re\alpha(\mathfrak p)=2a/t$,
\begin{equation}
 \left|\frac{e^{-(r+ih_*)^2/t}\mathcal M(w+r+ih_*)}
 {e^{-v_*^2/t}\mathcal M(\mathfrak p)}\right|
 \le e^{-d^2/t+CA}\le e^{-ct\mathcal L^2}.
 \label{eq:full-saddle-left-payment}
\end{equation}
Here $d\ge a/2$, $a\asymp t\mathcal L$ and $A=a+O(1)$;
the segment length divided by $\sqrt t$ is absorbed by reducing $c$.
Zeta is bounded by $\zeta(2)$ there.

At the top of each connector apply
\eqref{eq:full-saddle-left-integral} with $A-d=2$.
On $\Re=2$, $|\Im\alpha|$ is uniformly bounded, including bounded
height. Traversing the connector back toward $\tau_0$ costs at most
$e^{C|h_*|}$ in gamma modulus. Its Gaussian modulus cannot increase
relative to the top, because $\eta^2\le h_*^2$ along the segment.
Both connectors, including their lengths divided by $\sqrt t$, therefore
cost $Ce^{-ct\mathcal L^2}$ against their own carriers. A collapsed
connector at $h_*=0$ costs zero.

Euler summation and Stirling in bounded real strips, with exact
reflection to the left of $1/2$ and compactness at bounded height, give
some fixed finite $b_0$ such that
\begin{equation}
 \sup_{-1\le b\le2}|\xi(b+i\tau_0)/8|
 \le C(1+|\tau_0|)^{b_0}|\mathcal M(2+i\tau_0)|.
 \label{eq:full-saddle-middle}
\end{equation}
The right-hand factor never vanishes, so this estimate also covers
$\tau_0=0$. The endpoint comparison above bounds the full middle by
$C\mathcal R^{b_0}e^{-ct\mathcal L^2}$ relative to either carrier.
The middle and both infinite ends have thus been retained.

Finally, on $\Re(w+v)\ge2$, the complete ordinary Dirichlet series
\cite{DLMFZetaOuter2026} gives
$|\zeta(w+v)-1|\le3\,2^{-\Re(w+v)}$.
For $r=a+u\ge a/2$ its tilted Gaussian payment is
\begin{align}
 Ce^{-a\log2}\frac1{\sqrt{\pi t}}
       \int_\R e^{-u^2/(2t)-u\log2}\,du
 &\le Ce^{-a\log2+t(\log2)^2/2}\notag\\
 &\le C\mathcal R^{-t\log2/8}.
 \label{eq:full-saddle-arithmetic}
\end{align}
This uses $a\ge t\mathcal L/5$ after one onset uniform in $t\ge32$;
the other pieces already have \eqref{eq:full-saddle-left-payment}.
Every arithmetic mode is paid, without a finite-mode replacement.

Combining the complete pieces first gives the explicit three-part budget
\begin{equation}
 \frac{|H-P_*|}{|\mathcal S_t(s)|+|\mathcal S_t(1-s)|}
 \le C\left\{\frac{t^{3/2}}{N_*^2}
       +\mathcal R^{-t\log2/8}
       +\mathcal R^{b_0}e^{-ct\mathcal L^2}\right\}.
 \label{eq:full-saddle-three-payments}
\end{equation}
Because $t\log2/8>2$ at $t\ge32$ and $N_*\le C\mathcal R\log\mathcal R$,
the last two terms are absorbed by $Ce_*$ after an eventual onset.
This proves \eqref{eq:full-saddle-field}. The floor $32$ is an
arithmetic allowance, distinct from the regulator switch points $64,128$
and from any certified zero or asymptotic threshold.
\end{proof}

\subsection{Physical collars, potential jets and the regulator payment}
\begin{lemma}[Actual full-saddle collars]\label{lem:full-saddle-collars}
The field comparison supplies all the simple-root collars and
analytic-unit bounds needed by Lemma~\ref{lem:uniform-q-collar} on
the full averaging buffer. Moreover
\begin{equation}
 \|S_H-S_{P_*}\|_{L^1(\Omega)}
       +\|w_H-w_{P_*}\|_{L^1(\Omega)}\le Ce_*.
 \label{eq:full-saddle-L1}
\end{equation}
\end{lemma}
\begin{proof}
Differentiate the implicit formula itself. With every $\alpha$ evaluated
at $\mathfrak p=\mathfrak p_t(w)$,
\begin{align}
 (\mathfrak p_t)_w&=B^{-1},&
 D_*:=(\log\mathcal S_t)'(w)&=\alpha+\frac{t\alpha''}{4B^2},\notag\\
 D_*'&=\frac{\alpha'}B+\frac{t\alpha'''}{4B^3}
            +\frac{t^2(\alpha'')^2}{4B^4}=O(N_*^{-1}).
 \label{eq:full-saddle-spatial-derivatives}
\end{align}
Choose compatible logarithms and write
$P_*=2\mathcal A_*\cos\Theta_*$, where
\[
 \mathcal A_*=
  \exp\!\left(\frac{\log\mathcal S_t(s)+\log\mathcal S_t(1-s)}2\right),
 \qquad
 \Theta_*=
  \frac{\log\mathcal S_t(s)-\log\mathcal S_t(1-s)}{2i}.
\]
On real $x$ the carrier is positive and the phase is real. Directly,
\begin{align}
 \omega_*:=\partial_x\Theta_*(t,x)
 &=\tfrac12\Re D_*((1+ix)/2)\notag\\
 &=\tfrac14\log\!\left(\frac{|\mathfrak p_t((1+ix)/2)|}{2\pi}\right)
            +O(N_*^{-1}+tN_*^{-2})\asymp\mathcal L,\notag\\
 \Theta_{*,zz}&=O(N_*^{-1}),&
 |\partial_z\log\mathcal A_*|&\le C.
 \label{eq:full-saddle-phase-bounds}
\end{align}
The last bound retains cancellation of the large real logarithms in
$D_*(s)-D_*(1-s)$: on the real axis it equals $2i\Im D_*(s)$, which
is bounded, and \eqref{eq:full-saddle-spatial-derivatives} extends the
bound over the fixed transverse buffer. Thus the analytic carrier is
comparable there to $A_X=|\mathcal S_t((1+ix)/2)|$.

The phase coordinate is injective on disks of radius $c/\omega_*$.
The cosine lower bound and Rouch\'e's theorem, applied to
\eqref{eq:full-saddle-field}, give one real simple zero of each field
per disjoint collar, displacement $O(e_*/\omega_*)$, exterior floor
$cA_X$ and interior floor $cA_X\omega_*$ times distance to its actual
zero. Reality follows from conjugation and uniqueness. These collars
exhaust the zeros in the required strip; there are $O(\omega_*)$
meeting the support. The enlarged buffer contains the full disks of
edge collars. The same local phase argument and analytic-unit comparison
as in Lemma~\ref{lem:linear-collars} apply here, with
\eqref{eq:full-saddle-phase-bounds} and error $Ce_*$.
Lemma~\ref{lem:uniform-q-collar} now gives
\eqref{eq:full-saddle-L1} with the unchanged physical $q$, uniformly in
its size. The planar translation estimate, rather than a regulator
floor, pays the displaced roots.
\end{proof}

\begin{proof}[Completion of Theorem~\ref{thm:full-saddle-static}]
Transfer spatial derivatives of orders at most four to $K$ in
\eqref{eq:full-saddle-L1}. Its bounded weak derivatives have no boundary
delta at these orders. This proves the $U,W$ part of
\eqref{eq:full-saddle-static}, with the physical additive normalization;
no mass normalization or harmonic correction is made.

For density write exactly, as in Section~\ref{sec:protocol},
$\epsilon_g(t,u)=E_0(t)e^{-\beta(t)u^2/2}$, $r=\log q$.
Choose a sufficiently large fixed $B_0$ below.
If $\beta(t)X\le B_0$, the required derivatives
$r_x=-2\beta u$ and $\Delta r=-2\beta$ are uniformly bounded on the
full support, since $\beta$ itself is bounded on $t\ge32$.
Apply Lemma~\ref{lem:holomorphic-weak-density} to both fields, subtract,
and use \eqref{eq:full-saddle-L1}; also
$|w_H(1-w_H)-w_{P_*}(1-w_{P_*})|\le|w_H-w_{P_*}|$.
This gives the density bound including the unresolved regulator layer
and $t=64$.

For $\beta X>B_0$ put
\begin{equation}
 \lambda_X=
 \frac{\sup_{|u-x|\le\ell+2}\epsilon_g(t,u)}{A_X}.
 \label{eq:full-saddle-regulator-ratio}
\end{equation}
At growing time it is essential to retain the mass anchor, since a
separate $O(t\log^2t)$ estimate of $\log E_0$ is too costly at
$X\asymp t\log t$.
The complete-source normalization estimate of Appendix~\ref{app:source}
gives explicitly
\begin{equation}
 \left|\frac{H(t,0)}{M_0(t)}-1\right|
 \le\frac{D_0(0)}{\sqrt{2\pi}}=O(\kappa^{-1/4}).
 \label{eq:full-saddle-source-anchor}
\end{equation}
Indeed $H(t,0)=2\pi^2e^{q_t(L)}\kappa^{-1/2}\int_\R B_t(w)\,dw$
and $M_0=2\pi^2e^{q_t(L)}\sqrt{2\pi/\kappa}$; the appendix bounds
the difference of that integral from $\sqrt{2\pi}$ by $D_0(0)$.
All its complete-source gates hold eventually. The already proved field
estimate at $x=0$, where $P_*(t,0)=2\mathcal S_t(1/2)>0$, now gives
\begin{equation}
 cM_0(t)\le A_0(t)=\mathcal S_t(1/2)\le CM_0(t).
 \label{eq:full-saddle-mass-anchor}
\end{equation}
This is a consequence used in a bound; neither $H(t,0)$ nor a measured
mass defines the special-function model.

Integrate the bounded real-axis logarithmic carrier derivative in
\eqref{eq:full-saddle-phase-bounds} to obtain
$|\log(A_X/A_0)|\le C_1X$.
For $t\ge128$, $E_0=M_0/\kappa$ exactly, so for $X$ beyond the fixed
buffer,
\begin{equation}
 \log\lambda_X\le C-\log\kappa
       -\frac{(X-\ell-2)^2}{2\kappa}+C_1X.
 \label{eq:full-saddle-regulator-payment}
\end{equation}
Taking $B_0>8(C_1+1)$ makes $X/\kappa>B_0$ imply
$\lambda_X\le Ce^{-cX}$. The two-sided versions of
\eqref{eq:full-saddle-mass-anchor} and the amplitude bound also give
$|\log\lambda_X|\le C(1+X^2)$, using $\kappa<X/B_0$ and its eventual
positive lower bound. This retains the exact $M_0/\kappa$ payment.

The remaining times lie in a fixed compact interval $[32,T]$.
Sectorial expansion of \eqref{eq:full-saddle-model} gives
$\log A_X=-\pi X/8+O_T(\log^2X)$ there. Both $E_0$ and $\beta$ are
bounded. If $\beta X>B_0$, taking $B_0$ also larger than $\pi$ yields
the same exponential bound on $\lambda_X$, with the same polynomial
allowance for its logarithm. Small $X$ is excluded by enlarging $R_*$.
This covers the entire switch interval without assigning a positive
lower bound to $\beta$ immediately after its flat onset.

Apply the actual-collar regularization bound
\eqref{eq:linear-regulator-bias} separately to $F=H,P_*$ and its own
real zero measure $\nu_F$. It gives
\begin{equation}
 |R_F-K*\nu_F|\le C(1+X^2)\lambda_X^2
 \left[1+\frac{1+\log(1+C/\lambda_X^2)}{\omega_*}\right]
 \le Ce^{-cX}.
 \label{eq:full-saddle-regulator-bias}
\end{equation}
Both adaptive spatial curvature terms are included in this payment.
The $O(\omega_*)$ actual roots and their $O(e_*/\omega_*)$
displacement give $|K*\nu_H-K*\nu_{P_*}|\le Ce_*$.
In this regime either $t$ is bounded or $\kappa<X/B_0$, hence
$t\le C(1+X)$ and $N_*\le C(1+X)\log(2+X)$.
The exponential bias is therefore absorbed into $Ce_*$.
This completes the density proof. No zero-measure interpretation is
used in the complementary unresolved-regulator regime.
\end{proof}

\subsection{Uniform limits and the same-time inner overlap}
\begin{corollary}[Three static regimes]\label{cor:full-saddle-limits}
Each observable or spatial jet on the left of
\eqref{eq:full-saddle-static} has the following uniform bounds:
\begin{align}
 \sup_{x\in\R,\ |y|\le Y}\text{error}(t,x,y)
 &\le\frac C{\sqrt t\log^2(2+t)},&&t\ge T_*,
 \label{eq:full-saddle-time-limit}\\
 \sup_{t\ge32,\ |y|\le Y}\text{error}(t,x,y)
 &\le\frac C{\sqrt X\log^{3/2}(2+X)},&&X\ge X_*,
 \label{eq:full-saddle-height-limit}\\
 \sup_{32\le t\le T,\ |y|\le Y}\text{error}(t,x,y)
 &=O_T(X^{-2}),&&X\longrightarrow\infty.
 \label{eq:full-saddle-compact-limit}
\end{align}
The thresholds exist but are not numerically instantiated. In particular,
the outer threshold $X_*$ may be chosen independent of $t\ge32$,
and at late time the theorem includes every real height.
\end{corollary}
\begin{proof}
Use $N_*\ge t\log(2+t)$ for the first claim. For the second put
$u=t\log(2+X)/(1+X)$; boundedness of $u^{3/2}/(1+u)^2$ gives
\[
 \frac{t^{3/2}}{N_*^2}
 \le\frac{C}{\sqrt{1+X}\log^{3/2}(2+X)}
       \frac{u^{3/2}}{(1+u)^2}.
\]
For fixed $t\in[32,T]$, $N_*\ge1+X$ gives the third.
The large-time domain condition is satisfied at small $x$ by $t$ itself.
\end{proof}

\begin{corollary}[Actual static inner matching]\label{cor:inner-static-overlap}
Retain every source, carrier and collar gate and the full buffered inner
window of Theorem~\ref{thm:explicit}, with
$X_{\rm in}(t)\asymp\kappa^{1/3}$.
At sufficiently large time both approximations use exactly the same
physical $q$, and
\begin{equation}
 \|R_{P_*}-R_L\|_{\rm inner}
 \le C\kappa^{-1/2}+Ct^{-1/2}\log^{-2}(2+t)
 =O(\kappa^{-1/2}).
 \label{eq:full-saddle-inner-overlap}
\end{equation}
The eventual region $2X_*\le|x|\le X_{\rm in}(t)/2$ is nonempty
and leaves room for both buffers.
\end{corollary}
\begin{proof}
For $t\ge128$ the global regulator is exactly
\eqref{eq:positiveRegulator}. Apply the triangle inequality through
$R_H$, using \eqref{eq:explicitDensityRate} and
\eqref{eq:full-saddle-time-limit}. Since $\kappa\asymp t\log t$,
the second error is bounded by a constant times $\kappa^{-1/2}$.
The inner gates hold eventually and $X_{\rm in}(t)\to\infty$.
No termwise matching of unrelated saddle expansions is needed.
\end{proof}
This closes the static positive spatial gap at sufficiently large time.
The threshold $e^{4t}$ belongs to the earlier $P_{\rm lin}$ proof;
it is not a condition on the different field $P_*$.
Same-time spatial overlap does not provide overlap through finite heat
times, nor differentiated inner matching.

\subsection{The exact nonzero heat residual}
The implicit derivatives also give, directly,
\begin{equation}
 \partial_t\mathfrak p_t(w)=\frac\alpha{2B},\qquad
 \partial_t\log\mathcal S_t(w)
 =\frac{\alpha^2}4+\frac{\alpha'}{4B}
                         +\frac{t\alpha\alpha''}{8B^2}.
 \label{eq:full-saddle-time-derivatives}
\end{equation}
Combine these with \eqref{eq:full-saddle-spatial-derivatives} and
$\partial_{zz}\mathcal S_t((1\pm iz)/2)
=-\partial_{ww}\mathcal S_t/4$.
In $\partial_t\log\mathcal S_t-(D_*^2+D_*')/4$ the
$\alpha^2$, $\alpha'/B$ and $t\alpha\alpha''$ terms cancel exactly.
The residual is
\begin{equation}
 \boxed{\mathcal E_*:=(\partial_t+\partial_{zz})P_*
 =-\sum_{w=s,1-s}\mathcal S_t(w)
 \left\{\frac{t\alpha'''(\mathfrak p_t(w))}{16B^3}
  +\frac{5t^2[\alpha''(\mathfrak p_t(w))]^2}{64B^4}\right\}.}
 \label{eq:full-saddle-residual}
\end{equation}
Each $B$ in the sum is evaluated at its own branch
$\mathfrak p_t(w)$. In particular the residual is nonzero, with bound
\begin{equation}
 |\mathcal E_*|\le C\frac{t}{N_*^3}
       \bigl(|\mathcal S_t(s)|+|\mathcal S_t(1-s)|\bigr).
 \label{eq:full-saddle-residual-bound}
\end{equation}
Neither this field residual nor a derivative of a static error supplies
an evolution comparison by itself. Section~\ref{sec:positive-evolution}
proves the positive current, direct time, forcing/source and actual
covariance comparison using additional pole, weak-transfer and signed-stress
estimates. Its model equations retain the exact residual above.
The generic collar transfer excludes no collision-control family;
the canonical application additionally uses the unchanged xi completion,
reflection and full ordinary Dirichlet arithmetic. No RH or earlier-time
simplicity conclusion follows.
